\chapter*{Abstract}

Cube satellite technology has increased the availability of remote earth observation data. The \acrfull{hypso} project at the \acrfull{NTNU} employs hyperspectral imagers with 120 spectral bands to capture remote-sensing data. To facilitate effective storage and downlink of the large amounts of data produced, on-board image compression must be employed. The \acrfull{CCSDS} develops compression algorithms tailored for hyperspectral and multispectral images. Of these, the CCSDS 123.0-B-2 algorithm provides near-lossless compression, capable of high compression ratios. Intersample data dependencies pose challenges for the implementation of high-throughput hardware accelerators of the algorithm. This work evaluates the use of a novel weight update scheme, proposed by \citeauthor{jiaRemovalFeedbackLoop2025}, to increase the throughput of the accelerator developed for \acrshort{hypso} by \citeauthor{vorhaugHyperspectralImageCompression2024}. For verification of compressed images, an open-source decompressor is developed, based on the verification model written by \citeauthor{vorhaugHyperspectralImageCompression2024}. For \acrshort{hypso}-2 images, the novel weight update scheme achieves a minimal reduction in compression performance at under $1\%$. The modified accelerator reaches a total throughput of 202 MS/s on the Zynq ZU7EV \acrshort{FPGA}, with the smallest area footprint of known implementations at 6349 LUTs and 3480 FFs. A throughput of 351 MS/s is reached for the predictor module, making it the fastest in the open literature. The accelerator is successfully tested on-board the \acrshort{hypso}-2 satellite, running at a clock frequency of 92 MHz on the Zynq-7030.
